% ======================================================================
\chapter[Волновая функция. Волна де Бройля. Интерпретация Борна]%
        {Волновая функция. Свободная частица. Вероятностная интерпретация}
\label{t:04}
% ======================================================================

\begin{qbox}
\textbf{Билет 4.} Волновая функция. Волновая функция свободной частицы (волна де Бройля).
Вероятностная интерпретация волновой функции по М.~Борну.
\end{qbox}

\section{Содержание билета}

\textbf{Подтемы:}
\begin{enumerate}
 \item понятие волновой функции (состояния) и требования к ней; принцип суперпозиции;
 \item волновая функция свободной частицы: плоская волна де Бройля; фазовая и групповая скорости;
 \item волновой пакет, его движение и расплывание;
 \item вероятностная интерпретация Борна: плотность вероятности, нормировка, средние значения,
       распределение по импульсам; плотность потока вероятности.
\end{enumerate}

\textbf{Необходимые понятия:} гипотеза де Бройля (\tref{t:02}); комплексные числа,
плоская волна $\ee^{\ii(\vect k\vect r-\omega t)}$; интеграл Фурье; плотность распределения
случайной величины.

\section{Определения и физические принципы}

\begin{definition}[волновая функция]
Состояние квантовой частицы в момент $t$ полностью задаётся комплексной функцией координат
$\Psi(\vect r,t)$ — \emph{волновой функцией} (амплитудой вероятности). Всё, что можно узнать о
частице (вероятности результатов любых измерений), вычисляется по $\Psi$.
\end{definition}

\begin{principle}{Вероятностная интерпретация (М.~Борн, 1926)}
Величина
\begin{equation}\label{eq:4:born}
  \dd w=|\Psi(\vect r,t)|^2\,\dd V
\end{equation}
есть вероятность обнаружить частицу в момент $t$ в элементе объёма $\dd V$ около точки $\vect r$.
Функция $\rho(\vect r,t)=|\Psi|^2=\Psi^*\Psi$ — \emph{плотность вероятности}.
\end{principle}

\begin{principle}{Принцип суперпозиции}
Если $\Psi_1$ и $\Psi_2$ — возможные состояния системы, то возможно и состояние
$\Psi=c_1\Psi_1+c_2\Psi_2$ с любыми комплексными $c_1,c_2$. Вероятности при этом не
складываются: $|\Psi|^2=|c_1\Psi_1|^2+|c_2\Psi_2|^2+2\operatorname{Re}(c_1c_2^*\Psi_1\Psi_2^*)$ —
последний член описывает интерференцию.
\end{principle}

\emph{Физический смысл.} Волновая функция — не наблюдаемая величина, а амплитуда; наблюдаемы
лишь квадраты модулей (вероятности) и, косвенно, относительные фазы (через интерференцию).
Умножение $\Psi$ на постоянный фазовый множитель $\ee^{\ii\alpha}$ не меняет физического
состояния. Вероятности относятся к статистическому ансамблю одинаково приготовленных систем:
отдельное измерение даёт случайный результат.

\begin{definition}[стандартные условия]
Физически допустимая волновая функция должна быть: (1) однозначной; (2) непрерывной;
(3) конечной (ограниченной); (4) иметь непрерывные первые производные везде, где
потенциал конечен; (5) для связанных состояний — квадратично интегрируемой и нормированной:
\begin{equation}\label{eq:4:norm}
  \int|\Psi(\vect r,t)|^2\,\dd V=1 .
\end{equation}
\end{definition}

\emph{Область применимости.} Описание волновой функцией $\Psi(\vect r,t)$ относится к
нерелятивистской частице (без учёта спина; спин добавляет дискретную переменную, \tref{t:12}).
Для системы $N$ частиц $\Psi(\vect r_1,\dots,\vect r_N,t)$ — функция в $3N$-мерном
конфигурационном пространстве, а не волна в обычном пространстве.

\section{Основная теория}

\subsection{Почему нужна волновая функция}

Дифракция электронов (\tref{t:02}) показывает, что распределение частиц на экране описывается
законами интерференции волн, а каждая отдельная частица регистрируется в одной точке. Если
закрыть одну из двух щелей, картина меняется не на «половину», а качественно: распределение
при двух открытых щелях не равно сумме распределений от каждой щели. Значит, складываются не
вероятности, а амплитуды. Эту роль выполняет волновая функция: $\Psi=\Psi_1+\Psi_2$, а
вероятность — $|\Psi_1+\Psi_2|^2$. В опытах А.~Тономуры (1989) с поштучным прохождением
электронов через электронную бипризму интерференционная картина складывалась из отдельных
точек — наглядная демонстрация интерпретации Борна.

\subsection{Волновая функция свободной частицы}

Частица с определённым импульсом $\vect p$ и энергией $E=p^2/2m$ описывается плоской волной
де Бройля:
\begin{equation}\label{eq:4:plane}
  \Psi_{\vect p}(\vect r,t)=A\exp\Bigl[\frac{\ii}{\hbar}(\vect p\vect r-Et)\Bigr]
  =A\,\ee^{\ii(\vect k\vect r-\omega t)},\qquad
  \vect k=\frac{\vect p}{\hbar},\quad \omega=\frac{E}{\hbar}=\frac{\hbar k^2}{2m}.
\end{equation}
Выбор комплексной экспоненты, а не $\cos(\vect k\vect r-\omega t)$, существенен: для неё
$|\Psi|^2=|A|^2$ не зависит ни от $\vect r$, ни от $t$, что соответствует однородному
движению (частица равновероятно находится в любой точке), тогда как $\cos^2$ давал бы
«стоячие» узлы вероятности. Кроме того, именно экспонента — собственная функция оператора
импульса (\tref{t:05}).

\paragraph{Закон дисперсии и скорости.} Из $\omega(k)=\hbar k^2/2m$:
\begin{equation}\label{eq:4:vel}
  v_{\text{ф}}=\frac{\omega}{k}=\frac{\hbar k}{2m}=\frac{v}{2},\qquad
  v_{\text{гр}}=\frac{\dd\omega}{\dd k}=\frac{\hbar k}{m}=\frac{p}{m}=v .
\end{equation}
Фазовая скорость не имеет прямого физического смысла (она зависит даже от выбора начала
отсчёта энергии: прибавление $mc^2$ к $E$ меняет $\omega$, но не $\dd\omega/\dd k$). Скорость
частицы равна групповой скорости волнового пакета.

\paragraph{Нормировка плоской волны.} Интеграл $\int|\Psi_{\vect p}|^2\dd V=|A|^2\int\dd V$
расходится: плоская волна — идеализация (частица, которая с равной вероятностью может
находиться где угодно во всём бесконечном пространстве). Применяют два способа.
\begin{enumerate}
 \item \emph{Нормировка в ящике}: частица заключена в куб объёма $V=L^3$ с периодическими
       граничными условиями; тогда $A=1/\sqrt V$, а $\vect k=2\pi\vect n/L$ ($\vect n$ —
       вектор с целыми компонентами).
 \item \emph{Нормировка на $\delta$-функцию}: $A=(2\pi\hbar)^{-3/2}$, тогда
       $\int\Psi_{\vect p'}^*\Psi_{\vect p}\,\dd V=\delta(\vect p-\vect p')$ (вывод~\ref{der:4:delta}).
\end{enumerate}
Плоскую волну можно также трактовать как описание пучка частиц: при $|A|^2=n$ это пучок с
плотностью $n$ частиц на единицу объёма и плотностью потока $j=nv$ (см. ниже).

\subsection{Волновой пакет}

Реальная частица локализована в конечной области, и её состояние — суперпозиция плоских волн:
\begin{equation}\label{eq:4:packet}
  \Psi(x,t)=\frac{1}{\sqrt{2\pi}}\int_{-\infty}^{\infty}a(k)\,\ee^{\ii(kx-\omega(k)t)}\,\dd k ,
\end{equation}
где амплитуда $a(k)$ сосредоточена около $k_0$ в интервале ширины $\Delta k$. Пакет движется
как целое с групповой скоростью $v_{\text{гр}}=\hbar k_0/m$ (вывод~\ref{der:4:group}), т.\,е. со
скоростью классической частицы. Поскольку $\omega(k)$ нелинейна, разные компоненты движутся с
разными скоростями, и пакет \emph{расплывается}. Для гауссова пакета начальной ширины $\sigma_0$
(без вывода)
\begin{equation}\label{eq:4:spread}
  \sigma(t)=\sigma_0\sqrt{1+\Bigl(\frac{\hbar t}{2m\sigma_0^2}\Bigr)^2}.
\end{equation}
При больших $t$ $\sigma\approx\hbar t/(2m\sigma_0)=\Delta v\,t$, где $\Delta v=\Delta p/m=\hbar/(2m\sigma_0)$ —
разброс скоростей, т.\,е. расплывание — следствие соотношения неопределённостей. Оценка:
электрон, локализованный в области $\sigma_0=1$~Å, удваивает ширину пакета за
$t\sim2m\sigma_0^2/\hbar\approx2\cdot10^{-16}$~с; пылинка массой $10^{-12}$~г при
$\sigma_0=1$~мкм — за $\sim2\cdot10^{7}$~с, т.\,е. практически никогда.

\begin{figure}[htbp]
\centering
\begin{tikzpicture}
\begin{axis}[width=0.92\linewidth, height=4.6cm, axis lines=middle,
  xlabel={$x$}, ylabel={}, xmin=-7, xmax=13, ymin=-1.15, ymax=1.3,
  xtick=\empty, ytick=\empty, samples=400,
  xlabel style={at={(ticklabel* cs:1)}, anchor=west}]
  \addplot[domain=-6.5:6.5, accent, thin] {exp(-x^2/4)*cos(deg(5*x))};
  \addplot[domain=-6.5:6.5, accent, thick, dashed] {exp(-x^2/4)};
  \addplot[domain=-6.5:6.5, accent, thick, dashed] {-exp(-x^2/4)};
  \addplot[domain=3:13, keycol, thick] {0.62*exp(-(x-8)^2/(4*2.6^2/1))};
  \addplot[domain=3:13, keycol, thick, dashed] {-0.62*exp(-(x-8)^2/(4*2.6^2/1))};
  \node[font=\small, accent] at (axis cs:-3.6,1.1) {$t=0$: $\operatorname{Re}\Psi$ и $\pm|\Psi|$};
  \node[font=\small, keycol] at (axis cs:9.6,0.9) {$t>0$: огибающая};
  \draw[->, thick] (axis cs:1.5,1.15) -- (axis cs:5.5,1.15) node[midway, above, font=\small] {$v_{\text{гр}}$};
\end{axis}
\end{tikzpicture}
\caption{Гауссов волновой пакет: при $t=0$ показаны действительная часть $\Psi$ и огибающая
$\pm|\Psi|$; в более поздний момент пакет смещён на $v_{\text{гр}}t$ и расплылся
(амплитуда уменьшилась, ширина выросла; качественно).}
\label{fig:4:packet}
\end{figure}

\subsection{Вероятностная интерпретация: следствия}

\paragraph{Средние значения координаты.} Поскольку $|\Psi|^2$ — плотность вероятности,
\begin{equation}
  \avg{\vect r}=\int\vect r\,|\Psi|^2\dd V,\qquad \avg{f(\vect r)}=\int f(\vect r)|\Psi|^2\dd V .
\end{equation}

\paragraph{Распределение по импульсам.} Разложим $\Psi$ по плоским волнам (трёхмерное
обобщение~\eqref{eq:4:packet}):
\begin{equation}\label{eq:4:ap}
  \Psi(\vect r,t)=\int a(\vect p,t)\,\frac{\ee^{\ii\vect p\vect r/\hbar}}{(2\pi\hbar)^{3/2}}\,\dd^3p,
  \qquad
  a(\vect p,t)=\int\Psi(\vect r,t)\,\frac{\ee^{-\ii\vect p\vect r/\hbar}}{(2\pi\hbar)^{3/2}}\,\dd V .
\end{equation}
Согласно Борну, $|a(\vect p)|^2\dd^3p$ — вероятность того, что измерение импульса даст значение
в элементе $\dd^3p$ около $\vect p$. Равенство Парсеваля $\int|a|^2\dd^3p=\int|\Psi|^2\dd V=1$
обеспечивает согласованность нормировок. Это частный случай общего правила: коэффициенты
разложения $\Psi$ по собственным функциям оператора величины дают амплитуды вероятностей её
значений (\tref{t:05}).

\paragraph{Сохранение вероятности.} Вероятность найти частицу где-либо должна оставаться
равной 1. Из уравнения Шредингера (\tref{t:06}) следует локальный закон сохранения
\begin{equation}\label{eq:4:cont}
  \pdiff{\rho}{t}+\operatorname{div}\vect j=0,\qquad
  \vect j=\frac{\ii\hbar}{2m}\bigl(\Psi\nabla\Psi^*-\Psi^*\nabla\Psi\bigr)
  =\frac{\hbar}{m}\operatorname{Im}(\Psi^*\nabla\Psi),
\end{equation}
где $\vect j$ — \emph{плотность потока вероятности}. Для плоской волны~\eqref{eq:4:plane}
$\vect j=|A|^2\hbar\vect k/m=|A|^2\vect v$: поток = плотность $\times$ скорость, как для
классического пучка. Эта формула используется при вычислении коэффициентов отражения и
прохождения (\tref{t:07}).

\section{Математические выводы}

\begin{derivation}[label=der:4:group]{скорость движения волнового пакета}
\emph{Допущение:} $a(k)$ заметно отлична от нуля лишь в узкой окрестности $k_0$ ($\Delta k\ll k_0$),
и на рассматриваемых временах можно ограничиться линейным членом разложения
$\omega(k)\approx\omega_0+v_{\text{гр}}(k-k_0)$, где $\omega_0=\omega(k_0)$, $v_{\text{гр}}=(\dd\omega/\dd k)_{k_0}$.

Подставим в~\eqref{eq:4:packet} и вынесем множитель, не зависящий от $k$:
\begin{align*}
  \Psi(x,t)&\approx\frac{1}{\sqrt{2\pi}}\int a(k)\exp\bigl[\ii kx-\ii\omega_0t-\ii v_{\text{гр}}(k-k_0)t\bigr]\dd k\\
  &=\ee^{\ii(k_0v_{\text{гр}}-\omega_0)t}\cdot\frac{1}{\sqrt{2\pi}}\int a(k)\,\ee^{\ii k(x-v_{\text{гр}}t)}\dd k
  =\ee^{\ii(k_0v_{\text{гр}}-\omega_0)t}\,\Psi(x-v_{\text{гр}}t,\,0).
\end{align*}
Следовательно, $|\Psi(x,t)|^2=|\Psi(x-v_{\text{гр}}t,0)|^2$: распределение вероятности переносится
без изменения формы со скоростью $v_{\text{гр}}=\dd\omega/\dd k=\hbar k_0/m=p_0/m$. Учёт квадратичного
члена $\frac12\omega''(k-k_0)^2$, $\omega''=\hbar/m$, приводит к расплыванию~\eqref{eq:4:spread}.
\end{derivation}

\begin{derivation}[label=der:4:delta]{нормировка плоских волн на $\delta$-функцию}
Используем представление $\delta$-функции $\int_{-\infty}^{\infty}\ee^{\ii qx}\dd x=2\pi\,\delta(q)$.
Для одномерных волн $\psi_p(x)=A\ee^{\ii px/\hbar}$:
\[
  \int_{-\infty}^{\infty}\psi_{p'}^*(x)\psi_p(x)\,\dd x=|A|^2\int\ee^{\ii(p-p')x/\hbar}\dd x
  =|A|^2\,2\pi\,\delta\Bigl(\frac{p-p'}{\hbar}\Bigr)=|A|^2\,2\pi\hbar\,\delta(p-p').
\]
Условие $\int\psi_{p'}^*\psi_p\dd x=\delta(p-p')$ даёт $A=(2\pi\hbar)^{-1/2}$; в трёх измерениях
$A=(2\pi\hbar)^{-3/2}$. Волны с разными импульсами ортогональны; «нормировка» означает
полноту системы в смысле разложения~\eqref{eq:4:ap}.
\end{derivation}

\begin{derivation}[label=der:4:box]{нормировка в ящике с периодическими условиями}
Потребуем $\psi(x+L)=\psi(x)$ для $\psi=A\ee^{\ii kx}$: $\ee^{\ii kL}=1$ $\Rightarrow$
$k=2\pi n/L$, $n=0,\pm1,\dots$ Нормировка:
$\int_0^L|A|^2\dd x=|A|^2L=1$ $\Rightarrow$ $A=1/\sqrt L$. В трёх измерениях
$\psi_{\vect k}=V^{-1/2}\ee^{\ii\vect k\vect r}$, $\vect k=\frac{2\pi}{L}(n_x,n_y,n_z)$.
Число состояний в элементе $\dd^3k$ равно $V\dd^3k/(2\pi)^3$ (на одно значение $\vect k$
приходится объём $(2\pi/L)^3$), или $V\dd^3p/(2\pi\hbar)^3$ — это используется в статистике
(газ Ферми, \tref{t:17}; фазовый объём в бета-распаде, \tref{t:20}). При $L\to\infty$
физические результаты не зависят от $L$.
\end{derivation}

\begin{derivation}[label=der:4:current]{плотность потока для плоской волны}
Для $\Psi=A\ee^{\ii(\vect k\vect r-\omega t)}$ имеем $\nabla\Psi=\ii\vect k\Psi$,
$\nabla\Psi^*=-\ii\vect k\Psi^*$. Подставляя в~\eqref{eq:4:cont}:
\[
  \vect j=\frac{\ii\hbar}{2m}\bigl(\Psi(-\ii\vect k)\Psi^*-\Psi^*(\ii\vect k)\Psi\bigr)
  =\frac{\ii\hbar}{2m}(-2\ii\vect k)|A|^2=\frac{\hbar\vect k}{m}|A|^2=\vect v\,|A|^2 .
\]
Размерность при $[|A|^2]=\text{см}^{-3}$: $[\vect j]=\text{см}^{-2}\text{с}^{-1}$ — число частиц через
единичную площадку в единицу времени.
\end{derivation}

\section{Итоговые формулы}

\keyformula{\dd w=|\Psi(\vect r,t)|^2\dd V,\qquad \int|\Psi|^2\dd V=1}
 {$\Psi$ — волновая функция, $\dd w$ — вероятность нахождения в объёме $\dd V$.}
 {нерелятивистская частица, связанное (нормируемое) состояние.}
 {вероятности положения частицы, средние $\avg{f(\vect r)}=\int f|\Psi|^2\dd V$.}

\keyformula{\Psi_{\vect p}=A\exp\Bigl[\frac{\ii}{\hbar}(\vect p\vect r-Et)\Bigr],\qquad E=\frac{p^2}{2m}}
 {$\vect p$ — импульс, $E$ — энергия свободной частицы; $A=V^{-1/2}$ или $(2\pi\hbar)^{-3/2}$.}
 {свободная нерелятивистская частица с точно заданным импульсом.}
 {базис для разложения произвольных состояний; описание пучков частиц.}

\keyformula{v_{\text{ф}}=\frac{\omega}{k}=\frac v2,\qquad v_{\text{гр}}=\frac{\dd\omega}{\dd k}=\frac{p}{m}=v}
 {$\omega=\hbar k^2/2m$ — закон дисперсии волн де Бройля.}
 {свободная нерелятивистская частица; узкий пакет $\Delta k\ll k_0$.}
 {пакет движется со скоростью частицы; фазовая скорость физического смысла не имеет.}

\keyformula{\pdiff{\rho}{t}+\operatorname{div}\vect j=0,\qquad \vect j=\frac{\hbar}{m}\operatorname{Im}(\Psi^*\nabla\Psi)}
 {$\rho=|\Psi|^2$; $\vect j$ — плотность потока вероятности.}
 {действительный потенциал (нет поглощения).}
 {сохранение вероятности; коэффициенты отражения и прохождения.}

\section{Физическая интерпретация}

\textbf{Статистический характер.} Квантовая механика предсказывает не результат отдельного
измерения, а распределения вероятностей. Повторив опыт с одинаково приготовленными частицами
много раз, получаем распределение $|\Psi|^2$; отдельная частица регистрируется в одной точке.

\textbf{Волна и частица.} Интерференция электрона с самим собой (две щели) показывает, что
$\Psi$ описывает состояние одной частицы, а не ансамбля взаимодействующих частиц. Попытка
выяснить, через какую щель прошла частица, разрушает интерференцию: измерение меняет
состояние.

\textbf{Предельные случаи.} Плоская волна ($\Delta p=0$) полностью делокализована ($\Delta x=\infty$);
$\delta$-образный пакет ($\Delta x=0$) содержит все импульсы. Реальные состояния — промежуточные
пакеты, удовлетворяющие $\Delta x\Delta p\ge\hbar/2$. Для макрочастиц расплывание пренебрежимо, и
центр пакета движется по классической траектории.

\section{Что необходимо знать на экзамене}

\begin{examcheck}
\checkhead{Определения}
\begin{itemize}
 \item волновая функция, плотность вероятности, нормировка, стандартные условия;
 \item принцип суперпозиции и интерференционный член; фаза $\Psi$ не наблюдаема.
\end{itemize}
\checkhead{Формулы}
\begin{itemize}
 \item $\dd w=|\Psi|^2\dd V$; $\Psi_{\vect p}=A\,\ee^{\ii(\vect p\vect r-Et)/\hbar}$; $\omega=\hbar k^2/2m$;
 \item $v_{\text{ф}}=v/2$, $v_{\text{гр}}=v$; $\vect j=(\hbar/m)\operatorname{Im}(\Psi^*\nabla\Psi)$; $|a(\vect p)|^2$ — распределение по импульсам.
\end{itemize}
\checkhead{Выводы}
\begin{itemize}
 \item групповая скорость пакета равна скорости частицы;
 \item нормировка плоской волны (в ящике и на $\delta$-функцию); поток для плоской волны.
\end{itemize}
\checkhead{Интерпретации}
\begin{itemize}
 \item почему $\ee^{\ii kx}$, а не $\cos kx$; почему плоская волна не нормируема;
 \item опыт с двумя щелями и поштучной регистрацией; расплывание пакета и его оценки.
\end{itemize}
\end{examcheck}
